% 4.4 — Récupération et reçus
\section{La récupération et les reçus}
\label{sec:recuperation}

\subsection{Captures ciblées et groupes}
\fichier{recovery.py} copie la cible réelle avant l'exécution : un fichier, un répertoire
(avec ses sous-répertoires ignorés comme \fichier{.git} quand le parent est détruit), une base
SQLite. Depuis la partie S2, toutes les captures d'une commande partagent un identifiant de
\emph{groupe} : un seul \cmd{undo} restaure tout, et \cmd{snapshot\_set} est \emph{tout ou
rien} --- si une capture échoue, les précédentes sont invalidées et marquées comme telles dans
l'index. Les captures sont bornées (256\,Mo et \num{25000} fichiers par défaut,
réglables) ; au-delà, la capture est refusée et la commande escaladée, jamais réduite.

La restauration a ses propres pièges, découverts en les provoquant : un fichier rendu en
lecture seule par l'agent (\cmd{chmod 444} ou \cmd{000}) empêchait \cmd{undo} de le réécrire
(corrigé dans M2.2c) ; l'état écrasé par une restauration est lui-même sauvegardé, pour
qu'un \cmd{undo} malheureux se défasse à son tour.

\begin{terminal}[undo d'un groupe capturé par le filet]
$ demo_cli log
  id        when             kind     size    action
  8553dbf8  20261009-141447  group    30B     openat_trunc app.py (2 files)
  d9684862  20261009-141441  group    12B     rm -rf src (1 files)
$ demo_cli undo 8553dbf8
demo_cli 1.7.0  RESTORED
  target              ~/projet/src/app.py
  restored 2 files/databases as one undo
  target              ~/projet/notes.txt
  restored 2 files/databases as one undo
  the previous content was kept as edd72faa - put it back with:  demo_cli undo edd72faa
\end{terminal}

\subsection{L'aperçu SQL}
Pour une requête \cmd{DELETE}, \cmd{UPDATE} ou \cmd{TRUNCATE} sur une base locale,
\fichier{preview.py} compte en lecture seule les lignes qui seraient touchées (pilote standard
pour SQLite, session en lecture seule pour PostgreSQL). La décision devient DRY\_RUN :
capturée, réversible, avec le nombre de lignes affichées.

\subsection{Le registre de reçus}
Chaque décision --- y compris un ALLOW, et chaque capture faite par un filet --- laisse un reçu
dans \fichier{.demo\_cli/receipts.jsonl}. Un reçu contient la commande, la décision et sa
raison, la règle, la surface, le point de récupération, le contexte, l'intention déclarée, la
prédiction du moteur, le groupe, le fait qu'une question \emph{pouvait} être posée
(\cmd{approvable} --- jamais la réponse, qu'aucune porte ne voit) et l'invariant. Son
empreinte est calculée sur son contenu canonique \emph{et} l'empreinte du reçu précédent :

\begin{code}{Le chaînage (\fichier{receipts.py}, \cmd{Receipt.finalize})}
body = asdict(self)
body.pop("receipt_hash")
self.receipt_hash = hashlib.sha256(
    (_canon(body) + self.prev_receipt_hash).encode()).hexdigest()
\end{code}

Modifier, insérer, supprimer ou réordonner un reçu brise la chaîne. L'écriture est protégée par
un verrou multiplateforme (\cmd{fcntl} sous Linux, \cmd{msvcrt} sous Windows). Le filet Windows
écrit dans une seconde chaîne, liée à la première, pour éviter qu'un même fichier soit écrit
à travers le montage et directement. Les secrets (mots de passe d'URL, clés) sont masqués. Un
reçu réel complet figure en annexe~\ref{ann:recus}.

\begin{terminal}[vérification, puis après modification d'un seul champ]
$ demo_cli verify
demo_cli 1.7.0  VERIFIED
Receipt chains
  main         6 entries    INTACT
Recovery artifacts
  artifacts    4 verified on disk
$ # on change "decision" d'un reçu en ALLOW, puis :
$ demo_cli verify
demo_cli 1.7.0  TAMPERED
Receipt chains
  main         TAMPERED
             broken at line 2
             A field in this entry was edited after it was written.
\end{terminal}
